%%
% BIThesis 研究生学位论文模板 The BIThesis Template for Graduate Thesis
% This file has no copyright assigned and is placed in the Public Domain.
%%

\begin{abstract}
  大语言模型驱动的多智能体系统能够通过角色分工、协作推理和多轮通信提升复杂任务求解能力，但密集通信也会引入显著的上下文冗余、推理成本和工具调用开销。一方面，智能体之间传递的自然语言消息通常包含与接收方职责弱相关的推理过程、多个前驱节点的重复观点以及不断累积的历史信息，容易造成Token消耗增加和关键信息稀释。另一方面，在搜索、代码执行和文本检查等工具调用场景中，工具观察结果可能沿通信边在多个智能体之间传播，使单次工具调用开销被通信图结构进一步放大。针对上述问题，本文围绕多智能体系统中的高效通信技术展开研究，在不修改底层大语言模型参数的前提下，分别从通信内容和通信结构两个层面提出优化方法。

  首先，本文提出面向语义冗余的消息压缩框架。该框架将多智能体通信中的语义冗余归纳为角色不匹配、多前驱重复和跨轮历史累积三类，并分别设计角色感知压缩、节点感知聚合和分歧感知压缩模块。角色感知压缩根据接收方职责筛选发送方推理内容，节点感知聚合对多个前驱智能体的相似观点进行去重并保留来源信息，分歧感知压缩从多轮历史中提取共识、分歧及其简要依据，从而提高通信内容的信息密度。实验结果表明，该框架能够在多数知识问答、数学推理和代码生成任务上保持或提升任务准确率，并降低完全通信拓扑带来的额外推理成本。

  其次，本文提出面向通信边冗余的图剪枝算法。该方法将多智能体系统中的空间边和时间边建模为可学习概率，以任务正确性为收益项，并引入模型输出Token、工具调用次数和工具观察Token三类边级成本信号构建优化目标，使通信结构学习能够刻画不同发送方、不同轮次和不同工具观察带来的成本差异。在此基础上，本文引入Taylor一阶重要性评估，根据边参数及其梯度估计移除通信边对目标函数的影响，并周期性剪除低重要性边。实验结果表明，在AIME、GAIA、GPQA和BigCodeBench等工具调用任务中，该方法能够相对已有多智能体通信基线改善准确率--成本权衡，尤其在GAIA和GPQA等任务中表现出较好的通信边筛选能力。总体而言，本文从消息级压缩和图结构剪枝两个角度缓解多智能体系统中的语义冗余问题和通信边冗余问题，为构建低成本、可扩展的多智能体推理系统提供了可行方案。
\end{abstract}

% 如需手动控制换行连字符位置，可写 aa\-bb，详见
% https://bithesis.bitnp.net/faq/hyphen.html

\begin{abstractEn}
  Large-language-model-driven multi-agent systems improve complex task solving through role specialization, collaborative reasoning, and multi-round communication. However, dense communication also introduces substantial contextual redundancy, inference cost, and tool-use overhead. On the one hand, natural language messages exchanged among agents often contain reasoning traces weakly related to the receiver's responsibility, repeated opinions from multiple predecessors, and accumulated historical content, which increases token consumption and dilutes key information. On the other hand, in tool-augmented scenarios such as web search, code execution, and text checking, tool observations may propagate through multiple communication edges and further amplify the cost of a single tool call. To address these issues, this thesis studies efficient communication techniques for multi-agent systems and proposes two optimization methods from the perspectives of communication content and communication structure, without modifying the parameters of the underlying large language models.

  First, this thesis proposes a message compression framework for semantic redundancy. The framework categorizes semantic redundancy in multi-agent communication into three types: mismatch between sender messages and receiver roles, repeated information from multiple predecessor agents, and accumulation of multi-round historical context. Accordingly, it designs three modules, namely role-aware compression, node-aware aggregation, and disagreement-aware compression. Role-aware compression filters sender reasoning according to the responsibility of the receiver. Node-aware aggregation removes semantically similar statements from multiple predecessors while preserving source information. Disagreement-aware compression extracts consensus answers, conflicting candidates, and concise supporting rationales from historical interactions. Experimental results show that the proposed framework can maintain or improve task accuracy on most knowledge question answering, mathematical reasoning, and code generation tasks, while reducing the additional inference cost introduced by fully connected communication.

  Second, this thesis proposes a graph pruning algorithm for redundancy in communication edges. The method models spatial and temporal edges in a multi-agent system as learnable probabilities, and constructs an edge-level optimization objective by jointly considering task correctness, model output tokens, tool call frequency, and tool observation tokens. This design enables communication-structure learning to capture cost differences across senders, reasoning rounds, and tool observations. Furthermore, a first-order Taylor importance estimator is introduced to evaluate the influence of removing each edge based on its parameter value and gradient, and low-importance communication edges are pruned periodically. Experimental results on tool-augmented tasks including AIME, GAIA, GPQA, and BigCodeBench demonstrate that the proposed method improves the accuracy-cost trade-off over existing multi-agent communication baselines, especially on tasks that require evidence selection and tool-result integration. Overall, this thesis mitigates semantic redundancy and redundancy in communication edges through message compression and graph pruning, providing a practical approach for building cost-efficient and scalable multi-agent reasoning systems.
\end{abstractEn}
